\setcounter{subsection}{2}
\subsection{Quantum query and computational models}
For general background on the quantum query model we refer the reader to \cite{HLS07}. 
Here we restrict ourselves to describing the quantum 
implementation of the input oracles in the adjacency matrix and adjacency array models.

In the adjacency matrix model, on input a weighted graph $G = (V,w)$, classically one can query any $\{u,v\} \in V^{(2)}$ and receive the answer 
$w(\{u,v\})$.  We now describe how to model this by a quantum query.  We will assume that the 
edge weights are given as binary decimal numbers with $M_1$ bits before the decimal and $M_2$ bits after the decimal for a total of $M = M_1 + M_2$ bits.  
The state of the quantum query 
algorithm will have three registers, a query register, an answer register, and a workspace register.  The 
state of the algorithm will in general be in a superposition of the basis states $|\{u,v\} \rangle |b \rangle | a \rangle$ where $\{u,v\} \in V^{(2)}, b \in \{0,1\}^M$ 
and $a \in \Acal$ for an arbitrary finite set $\Acal$.  
On input graph $G = (V,w)$, the input oracle $\mathsf{O}_G$ acts on a basis state $|\{u,v\} \rangle |b \rangle | a \rangle$ as
\[
\mathsf{O}_G |\{u,v\} \rangle |b \rangle | a \rangle = |\{u,v\} \rangle |b \oplus w(\{u,v\}) \rangle | a \rangle \enspace.
\]
 
In the adjacency array model, on input a weighted $n$-vertex graph $G = (V,w)$ one can make two types of queries.  In the first type, one can query a vertex $v \in V$ and receive its degree $\deg(v)$.
The second type is specified by a family of functions $\{f_v : [\deg(v)] \rightarrow V\}_{v \in V}$ such that $f_v(i)$ corresponds to the $\ith$ neighbor of vertex $v$ (according to some arbitrary but fixed ordering).  
A query consists of a pair $(v,i)$ for $i \in [\deg(v)]$ and the returned answer is the pair
$(f_v(i), w(\{v,f_v(i)\}))$.  In this paper we will only 
need to model the second type of query quantumly.  This is because our upper bounds are larger than $n$ so we can let the algorithm classically 
query all degrees at the start of the algorithm, and in our lower bound on the query complexity of edge connectivity for weighted graphs we assume the 
algorithm already knows the degree of every vertex.  
The state of the quantum query algorithm will again have a query register, an answer register, and a workspace register, with the state of the algorithm 
in general being in a superposition of the basis states $|(v,i)\rangle |x \rangle |b \rangle | a \rangle$ where $v \in V, i \in [\deg(v)], x \in \{0,\ldots, n-1\}, b \in \{0,1\}^M$, 
and $a \in \Acal$ for an arbitrary finite set $\Acal$.  We further let $\tau : V \rightarrow \{0,1,\ldots, n-1\}$ be a bijection where $|V| = n$.  Then the input oracle $\mathsf{O}_G$ acts on a basis state in the following way:
\[
\mathsf{O}_G |(v,i)\rangle |x \rangle |b \rangle | a \rangle = |(v,i) \rangle |x + \tau(f_v(i)) \bmod n \rangle |b \oplus w(\{v,f_v(i)\}) \rangle | a \rangle \enspace.
\]

In Section~4 of the source we will further show that our query algorithms can be implemented in a time efficient manner.  We analyze the time complexity in terms of the standard quantum circuit model augmented with two types of oracles.  
One is the oracle for the input, either in the adjacency matrix or array model, and the second is an oracle to a classical memory of $\tO(n)$ bits.  The latter corresponds to a \emph{quantum random-access-memory} or \emph{QRAM}.  We further assume that 
we can classically update a value in this $\tO(n)$ bit classical memory in time $\tO(1)$.  
The assumption of QRAM access is also required for the time efficiency of the sparsifier construction in \cite{AdW19} which our algorithms build on, and in fact is a necessary (but sometimes inexplicit) assumption in the time analysis of many quantum algorithms for graph problems, e.g.~\cite{DHHM06,Amb06,B13}.

